\documentclass[a4paper,12pt]{article}

\usepackage[T2A]{fontenc}
\usepackage[utf8]{inputenc}
\usepackage[russian]{babel}
\usepackage{amsmath,amssymb,amsthm}
\usepackage{array}
\usepackage{booktabs}
\usepackage{enumitem}
\usepackage{listings}
\usepackage{xcolor}
\usepackage{geometry}
\usepackage{tikz}

\geometry{margin=2cm}

\newtheorem{definition}{Определение}
\newtheorem{theorem}{Теорема}
\newtheorem{statement}{Утверждение}

\lstset{
  basicstyle=\ttfamily\small,
  keywordstyle=\color{blue},
  commentstyle=\color{gray},
  frame=single,
  breaklines=true
}

\title{Билет №24 \\ \small Оперативная память. Способы адресации. Реальный и защищенный режим работы процессора. Виртуальная память. Страничная организация памяти. Файлы подкачки, алгоритмы выгрузки страниц. (ИТМО). \\ Адресация, адресные преобразования. Сегменты, страницы, регионы. Аппаратные средства поддержки виртуальной памяти. (СпБГУ)}
\author{Сериков Владислав Александрович}
\date{4 мая 2026}

\begin{document}

\maketitle
\tableofcontents
\newpage

\section{Иерархия памяти и оперативная память}

\subsection{Иерархия памяти}

Память вычислительной системы организована иерархически:

\[
\text{регистры} \rightarrow \text{кэш L1/L2/L3} \rightarrow
\text{оперативная память} \rightarrow \text{SSD/HDD} \rightarrow
\text{архивные носители}.
\]

Чем выше уровень иерархии, тем меньше объём, ниже задержка доступа и выше стоимость одного байта.

\begin{definition}
\textbf{Оперативная память} --- энергозависимая память с произвольным доступом, используемая процессором для хранения исполняемого кода, данных процессов, структур ядра ОС, буферов ввода-вывода и кэшей файловой системы.
\end{definition}

Оперативная память обычно реализуется на основе DRAM. Её важные свойства:

\begin{itemize}
    \item \textbf{произвольный доступ}: время доступа слабо зависит от адреса;
    \item \textbf{энергозависимость}: данные теряются при отключении питания;
    \item \textbf{байтовая адресация}: минимальная адресуемая единица обычно байт;
    \item \textbf{использование совместно с кэшами}: процессор чаще обращается к кэшу, а не напрямую к DRAM.
\end{itemize}

\subsection{Адрес, адресное пространство, слово памяти}

\begin{definition}
\textbf{Адрес} --- числовой идентификатор ячейки памяти или блока памяти.
\end{definition}

\begin{definition}
\textbf{Адресное пространство} --- множество адресов, доступных некоторому субъекту: процессору, процессу, устройству или операционной системе.
\end{definition}

Если адрес имеет $n$ бит, то максимальное число различных адресов равно

\[
2^n.
\]

При байтовой адресации объём адресуемой памяти равен

\[
2^n \text{ байт}.
\]

Например:

\[
32\text{-битный адрес} \Rightarrow 2^{32} \text{ байт} = 4 \text{ ГиБ}.
\]

\[
64\text{-битный адрес} \Rightarrow 2^{64} \text{ байт} = 16 \text{ ЭиБ}.
\]

На практике современные процессоры часто используют не все 64 бита виртуального или физического адреса.

\section{Способы адресации}

\subsection{Понятие способа адресации}

\begin{definition}
\textbf{Способ адресации} --- правило, по которому процессор вычисляет адрес операнда команды.
\end{definition}

Операнд команды может находиться:

\begin{itemize}
    \item непосредственно в самой команде;
    \item в регистре;
    \item в оперативной памяти;
    \item по адресу, вычисляемому из нескольких составляющих.
\end{itemize}

\subsection{Основные способы адресации}

\subsubsection{Непосредственная адресация}

Операнд записан прямо в инструкции.

\[
\texttt{MOV AX, 5}
\]

Здесь число $5$ является непосредственным операндом.

\subsubsection{Регистровая адресация}

Операнд находится в регистре процессора.

\[
\texttt{ADD AX, BX}
\]

Оба операнда находятся в регистрах.

\subsubsection{Прямая абсолютная адресация}

В команде указан адрес памяти.

\[
\texttt{MOV AX, [1000h]}
\]

Процессор читает данные из ячейки памяти с адресом $1000_h$.

\subsubsection{Косвенная регистровая адресация}

Адрес операнда хранится в регистре.

\[
\texttt{MOV AX, [BX]}
\]

Если $BX = 2000_h$, то обращение идёт к памяти по адресу $2000_h$.

\subsubsection{Базовая адресация}

Эффективный адрес вычисляется как сумма базового регистра и смещения:

\[
EA = Base + Displacement.
\]

Пример:

\[
\texttt{MOV AX, [BP + 8]}
\]

Такой способ часто используется для доступа к параметрам функции или локальным переменным.

\subsubsection{Индексная адресация}

Эффективный адрес вычисляется через индексный регистр:

\[
EA = Index + Displacement.
\]

Пример:

\[
\texttt{MOV AX, [SI + 4]}
\]

Используется для доступа к элементам массивов.

\subsubsection{Базово-индексная адресация}

\[
EA = Base + Index + Displacement.
\]

Пример:

\[
\texttt{MOV AX, [BX + SI + 10h]}
\]

\subsubsection{Масштабированная индексная адресация}

В архитектуре x86-64 часто используется форма

\[
EA = Base + Index \cdot Scale + Displacement,
\]

где

\[
Scale \in \{1, 2, 4, 8\}.
\]

Пример:

\[
\texttt{MOV EAX, [RBX + RCX * 4 + 16]}
\]

Если массив состоит из 4-байтовых элементов, то $RCX$ может хранить номер элемента.

\subsection{Минимальный пример вычисления эффективного адреса}

Пусть:

\[
RBX = 1000_h,\quad RCX = 3,\quad Scale = 4,\quad Displacement = 20_h.
\]

Тогда

\[
EA = 1000_h + 3 \cdot 4 + 20_h.
\]

Так как $3 \cdot 4 = 12 = C_h$, получаем:

\[
EA = 1000_h + C_h + 20_h = 102C_h.
\]

\section{Виды адресов и адресные преобразования}

\subsection{Логический, линейный, виртуальный и физический адрес}

В системах с сегментацией и страничной организацией памяти различают несколько видов адресов.

\begin{definition}
\textbf{Логический адрес} --- адрес, используемый программой до применения сегментного преобразования. В x86 логический адрес имеет вид:
\[
\text{selector} : \text{offset}.
\]
\end{definition}

\begin{definition}
\textbf{Линейный адрес} --- адрес, полученный после сегментного преобразования.
\end{definition}

\begin{definition}
\textbf{Виртуальный адрес} --- адрес, которым оперирует процесс в своём виртуальном адресном пространстве.
\end{definition}

В простых моделях виртуальный адрес часто отождествляют с линейным адресом.

\begin{definition}
\textbf{Физический адрес} --- адрес в реальной оперативной памяти, поступающий на шину памяти или в контроллер памяти.
\end{definition}

В общем случае схема преобразования такова:

\[
\text{логический адрес}
\rightarrow
\text{линейный адрес}
\rightarrow
\text{физический адрес}.
\]

Для x86 в защищённом режиме:

\[
(\text{selector}, \text{offset})
\rightarrow
\text{linear address}
\rightarrow
\text{physical address}.
\]

\subsection{Общая схема трансляции адреса}

\[
\boxed{
\text{Виртуальный адрес}
=
\text{номер страницы}
+
\text{смещение внутри страницы}
}
\]

\[
\boxed{
\text{Физический адрес}
=
\text{номер кадра}
+
\text{смещение внутри кадра}
}
\]

Если размер страницы равен $2^k$ байт, то младшие $k$ бит адреса задают смещение, а старшие биты --- номер страницы.

\[
VA = VPN \cdot 2^k + offset,
\]

где:

\begin{itemize}
    \item $VA$ --- виртуальный адрес;
    \item $VPN$ --- virtual page number, номер виртуальной страницы;
    \item $offset$ --- смещение внутри страницы;
    \item $2^k$ --- размер страницы.
\end{itemize}

После обращения к таблице страниц:

\[
VPN \mapsto PFN,
\]

где $PFN$ --- physical frame number.

Физический адрес:

\[
PA = PFN \cdot 2^k + offset.
\]

Важно, что смещение внутри страницы при трансляции не меняется.

\section{Реальный режим работы процессора}

\subsection{Общая характеристика}

\begin{definition}
\textbf{Реальный режим} --- режим работы x86-процессора, совместимый с Intel 8086, в котором используется сегментная адресация вида:
\[
\text{segment} : \text{offset}.
\]
\end{definition}

В классическом реальном режиме физический адрес вычисляется так:

\[
PA = 16 \cdot Segment + Offset.
\]

Иначе:

\[
PA = (Segment \ll 4) + Offset.
\]

Здесь:

\begin{itemize}
    \item $Segment$ --- значение сегментного регистра;
    \item $Offset$ --- смещение внутри сегмента;
    \item $PA$ --- физический адрес.
\end{itemize}

\subsection{Пример вычисления адреса в реальном режиме}

Пусть:

\[
CS = 1234_h,\quad IP = 5678_h.
\]

Тогда физический адрес инструкции:

\[
PA = 1234_h \cdot 10_h + 5678_h.
\]

\[
1234_h \cdot 10_h = 12340_h.
\]

\[
PA = 12340_h + 5678_h = 179B8_h.
\]

\subsection{Свойства реального режима}

\begin{itemize}
    \item адресация через сегмент и смещение;
    \item отсутствие полноценной защиты памяти;
    \item отсутствие изоляции процессов;
    \item отсутствие аппаратной виртуальной памяти в современном смысле;
    \item возможность обращения примерно к первому мегабайту памяти в классической модели;
    \item используется при начальной загрузке компьютера, затем ОС обычно переводит процессор в защищённый или long mode.
\end{itemize}

\subsection{Перекрытие сегментов}

В реальном режиме разные пары segment:offset могут указывать на один и тот же физический адрес.

Например:

\[
1000_h:0010_h
\]

даёт:

\[
1000_h \cdot 10_h + 0010_h = 10010_h.
\]

А

\[
1001_h:0000_h
\]

даёт:

\[
1001_h \cdot 10_h + 0000_h = 10010_h.
\]

Следовательно:

\[
1000_h:0010_h = 1001_h:0000_h.
\]

\section{Защищённый режим работы процессора}

\subsection{Общая характеристика}

\begin{definition}
\textbf{Защищённый режим} --- режим x86-процессора, в котором поддерживаются аппаратные механизмы защиты памяти, сегментация, страничная организация, уровни привилегий и изоляция задач.
\end{definition}

Основные возможности защищённого режима:

\begin{itemize}
    \item расширенное адресное пространство;
    \item аппаратная защита памяти;
    \item сегменты с дескрипторами;
    \item таблицы дескрипторов GDT и LDT;
    \item уровни привилегий;
    \item страничная организация памяти;
    \item виртуальная память;
    \item многозадачность с изоляцией процессов.
\end{itemize}

\subsection{Сегментная модель в защищённом режиме}

В защищённом режиме сегментный регистр содержит не сам адрес сегмента, а \textbf{селектор сегмента}.

\begin{definition}
\textbf{Селектор сегмента} --- значение, указывающее на дескриптор сегмента в таблице дескрипторов.
\end{definition}

Селектор содержит:

\begin{itemize}
    \item индекс дескриптора;
    \item бит выбора таблицы: GDT или LDT;
    \item уровень привилегий запроса.
\end{itemize}

\begin{definition}
\textbf{Дескриптор сегмента} --- структура, описывающая сегмент: базовый адрес, предел, права доступа, тип, уровень привилегий и служебные признаки.
\end{definition}

Логический адрес имеет вид:

\[
\text{selector} : \text{offset}.
\]

Сначала по селектору находится дескриптор. Из дескриптора извлекаются:

\[
Base,\quad Limit,\quad Rights.
\]

Затем проверяется условие:

\[
0 \leq Offset \leq Limit.
\]

Если проверка успешна, вычисляется линейный адрес:

\[
LA = Base + Offset.
\]

Если включена страничная организация, то линейный адрес дополнительно преобразуется в физический.

\subsection{Уровни привилегий}

В x86 используются четыре кольца защиты:

\[
0,1,2,3.
\]

\begin{itemize}
    \item Ring 0 --- максимальные привилегии, обычно ядро ОС;
    \item Ring 3 --- минимальные привилегии, пользовательские программы;
    \item Ring 1 и Ring 2 используются редко.
\end{itemize}

Цель уровней привилегий --- запретить пользовательским программам выполнять опасные операции напрямую.

\subsection{Сравнение реального и защищённого режимов}

\begin{center}
\begin{tabular}{|p{4cm}|p{5cm}|p{5cm}|}
\hline
\textbf{Свойство} & \textbf{Реальный режим} & \textbf{Защищённый режим} \\
\hline
Адресация & segment:offset, физический адрес $Segment \cdot 16 + Offset$ & selector:offset, через дескрипторы сегментов \\
\hline
Защита памяти & Практически отсутствует & Есть аппаратная защита \\
\hline
Изоляция процессов & Нет & Есть \\
\hline
Виртуальная память & Нет в современном смысле & Поддерживается через paging \\
\hline
Уровни привилегий & Нет & Есть кольца защиты \\
\hline
Типичное применение & Начальная загрузка, совместимость & Работа ОС, многозадачность \\
\hline
\end{tabular}
\end{center}

\section{Сегменты, страницы и регионы}

\subsection{Сегменты}

\begin{definition}
\textbf{Сегмент} --- логически связанная область адресного пространства, например код, данные, стек или область динамической памяти.
\end{definition}

Сегментация отражает логическую структуру программы.

Пример сегментов:

\begin{itemize}
    \item сегмент кода;
    \item сегмент данных;
    \item сегмент стека;
    \item сегмент библиотек;
    \item сегмент памяти ядра.
\end{itemize}

Преимущества сегментации:

\begin{itemize}
    \item естественное разделение кода, данных и стека;
    \item разные права доступа для разных областей;
    \item возможность совместного использования сегментов;
    \item защита от выхода за пределы сегмента.
\end{itemize}

Недостатки:

\begin{itemize}
    \item внешняя фрагментация;
    \item сложность выделения непрерывных областей;
    \item необходимость уплотнения памяти при сильной фрагментации.
\end{itemize}

\subsection{Страницы}

\begin{definition}
\textbf{Страница} --- блок виртуального адресного пространства фиксированного размера.
\end{definition}

\begin{definition}
\textbf{Кадр страницы} --- блок физической памяти того же размера, что и страница.
\end{definition}

Страничная организация разбивает виртуальную память на страницы, а физическую память --- на кадры.

Если размер страницы равен $P$, то:

\[
P = 2^k.
\]

Виртуальный адрес делится на:

\[
\text{номер страницы} \;|\; \text{смещение}.
\]

Физический адрес делится на:

\[
\text{номер кадра} \;|\; \text{смещение}.
\]

Преимущества страниц:

\begin{itemize}
    \item отсутствие внешней фрагментации;
    \item простое выделение памяти блоками фиксированного размера;
    \item возможность demand paging;
    \item удобная реализация виртуальной памяти;
    \item возможность совместного использования страниц между процессами.
\end{itemize}

Недостатки:

\begin{itemize}
    \item внутренняя фрагментация;
    \item необходимость таблиц страниц;
    \item накладные расходы на трансляцию адресов;
    \item page fault при отсутствии страницы в ОЗУ.
\end{itemize}

\subsection{Регионы}

\begin{definition}
\textbf{Регион памяти} --- непрерывная область виртуального адресного пространства процесса с одинаковыми или близкими свойствами: правами доступа, назначением, способом отображения.
\end{definition}

Примеры регионов:

\begin{itemize}
    \item регион исполняемого кода;
    \item регион глобальных данных;
    \item куча;
    \item стек;
    \item отображённый в память файл;
    \item разделяемая библиотека;
    \item память, выделенная через \texttt{mmap}.
\end{itemize}

Регион обычно описывается структурой вида:

\[
(start, end, permissions, backing\_object).
\]

Где:

\begin{itemize}
    \item $start$ --- начало региона;
    \item $end$ --- конец региона;
    \item $permissions$ --- права доступа;
    \item $backing\_object$ --- файл, swap или анонимная память.
\end{itemize}

\section{Виртуальная память}

\subsection{Определение и идея}

\begin{definition}
\textbf{Виртуальная память} --- механизм, при котором каждому процессу предоставляется собственное виртуальное адресное пространство, а аппаратно-программная система отображает виртуальные адреса в физические адреса или во внешнее хранилище.
\end{definition}

Основные цели виртуальной памяти:

\begin{itemize}
    \item изоляция процессов;
    \item защита памяти;
    \item упрощение программирования;
    \item возможность запускать программы, суммарный размер которых больше объёма ОЗУ;
    \item эффективное совместное использование памяти;
    \item поддержка отображения файлов в память;
    \item реализация copy-on-write.
\end{itemize}

\subsection{Иллюзия непрерывного адресного пространства}

Каждый процесс видит своё адресное пространство как большое и почти непрерывное:

\[
0x00000000 \ldots 0xFFFFFFFF
\]

для 32-битной модели или значительно большее для 64-битной.

Но физически страницы процесса могут быть размещены в разных кадрах оперативной памяти:

\[
\begin{array}{c}
\text{виртуальная страница 0} \rightarrow \text{физический кадр 5},\\
\text{виртуальная страница 1} \rightarrow \text{физический кадр 17},\\
\text{виртуальная страница 2} \rightarrow \text{отсутствует в ОЗУ},\\
\text{виртуальная страница 3} \rightarrow \text{физический кадр 2}.
\end{array}
\]

\subsection{Основные механизмы виртуальной памяти}

Виртуальная память строится на следующих механизмах:

\begin{enumerate}
    \item Разбиение виртуальной памяти на страницы.
    \item Разбиение физической памяти на кадры.
    \item Таблицы страниц.
    \item Аппаратный блок трансляции адресов MMU.
    \item Кэш трансляций TLB.
    \item Исключение page fault.
    \item Файл или раздел подкачки.
    \item Алгоритмы замещения страниц.
\end{enumerate}

\section{Страничная организация памяти}

\subsection{Таблица страниц}

\begin{definition}
\textbf{Таблица страниц} --- структура данных, отображающая номера виртуальных страниц в номера физических кадров и содержащая признаки доступа.
\end{definition}

Запись таблицы страниц обычно содержит:

\begin{itemize}
    \item номер физического кадра;
    \item бит присутствия;
    \item биты прав доступа;
    \item бит изменения;
    \item бит обращения;
    \item бит пользовательского/системного доступа;
    \item бит запрета исполнения;
    \item служебные признаки кэширования.
\end{itemize}

Обозначим запись таблицы страниц как PTE.

Типичная структура PTE:

\[
PTE = (PFN, P, R/W, U/S, A, D, NX).
\]

Где:

\begin{itemize}
    \item $PFN$ --- номер физического кадра;
    \item $P$ --- present bit;
    \item $R/W$ --- разрешение записи;
    \item $U/S$ --- пользовательская или системная страница;
    \item $A$ --- accessed bit;
    \item $D$ --- dirty bit;
    \item $NX$ --- запрет исполнения.
\end{itemize}

\subsection{Одноуровневая таблица страниц}

Пусть:

\begin{itemize}
    \item виртуальный адрес имеет 32 бита;
    \item размер страницы равен $4$ КиБ $= 2^{12}$ байт.
\end{itemize}

Тогда:

\[
VA = VPN \;|\; offset.
\]

Смещение занимает 12 бит.

\[
VPN \text{ занимает } 32 - 12 = 20 \text{ бит}.
\]

Количество виртуальных страниц:

\[
2^{20}.
\]

Если одна запись таблицы страниц занимает 4 байта, размер таблицы:

\[
2^{20} \cdot 4 = 4 \text{ МиБ}.
\]

Для одного процесса это ещё приемлемо, но для большого количества процессов становится дорого.

\subsection{Многоуровневые таблицы страниц}

Чтобы не хранить огромную непрерывную таблицу страниц, используют многоуровневые таблицы.

Например, 32-битный адрес при двухуровневой схеме может делиться так:

\[
\underbrace{10}_{\text{каталог}}
\;|\;
\underbrace{10}_{\text{таблица}}
\;|\;
\underbrace{12}_{\text{смещение}}.
\]

Трансляция:

\[
VA = (p_1, p_2, offset).
\]

\begin{enumerate}
    \item По $p_1$ выбирается запись каталога страниц.
    \item Эта запись указывает на таблицу страниц второго уровня.
    \item По $p_2$ выбирается PTE.
    \item Из PTE получается номер физического кадра.
    \item К нему присоединяется смещение.
\end{enumerate}

Физический адрес:

\[
PA = PFN \cdot 2^{12} + offset.
\]

\subsection{Пример трансляции адреса}

Пусть:

\begin{itemize}
    \item виртуальный адрес: $VA = 0x12345678$;
    \item размер страницы: $4$ КиБ;
    \item используется двухуровневая схема $10|10|12$.
\end{itemize}

Смещение:

\[
offset = VA \& 0xFFF = 0x678.
\]

Индекс таблицы второго уровня:

\[
p_2 = (VA >> 12) \& 0x3FF.
\]

Индекс каталога:

\[
p_1 = (VA >> 22) \& 0x3FF.
\]

После вычисления:

\[
p_1 = 72,\quad p_2 = 837,\quad offset = 0x678.
\]

Если PTE указывает на физический кадр

\[
PFN = 0xABC,
\]

то физический адрес:

\[
PA = 0xABC000 + 0x678 = 0xABC678.
\]

\section{Аппаратная поддержка виртуальной памяти}

\subsection{MMU}

\begin{definition}
\textbf{MMU} --- memory management unit, аппаратный блок, выполняющий трансляцию виртуальных адресов в физические и проверку прав доступа.
\end{definition}

MMU выполняет:

\begin{itemize}
    \item разбор виртуального адреса;
    \item обращение к таблицам страниц;
    \item проверку бита присутствия;
    \item проверку прав доступа;
    \item формирование физического адреса;
    \item генерацию исключения при ошибке.
\end{itemize}

\subsection{TLB}

\begin{definition}
\textbf{TLB} --- translation lookaside buffer, ассоциативный кэш трансляций виртуальных страниц в физические кадры.
\end{definition}

TLB хранит пары:

\[
(VPN, ASID) \mapsto (PFN, permissions).
\]

Здесь ASID --- идентификатор адресного пространства, позволяющий различать одинаковые виртуальные адреса разных процессов.

Если ASID не используется, при переключении процесса часто приходится очищать TLB или частично инвалидировать его.

\subsection{Алгоритм трансляции адреса с TLB}

\textbf{Вход:} виртуальный адрес $VA$.

\textbf{Выход:} физический адрес $PA$ или исключение.

\begin{enumerate}
    \item Разделить $VA$ на номер виртуальной страницы $VPN$ и смещение $offset$.
    \item Найти $VPN$ в TLB.
    \item Если запись найдена:
    \begin{enumerate}
        \item проверить права доступа;
        \item если права нарушены, сгенерировать исключение защиты;
        \item иначе получить $PFN$;
        \item сформировать $PA = PFN \cdot P + offset$.
    \end{enumerate}
    \item Если записи нет:
    \begin{enumerate}
        \item выполнить page table walk;
        \item если страница отсутствует, сгенерировать page fault;
        \item если страница присутствует, проверить права;
        \item загрузить трансляцию в TLB;
        \item сформировать физический адрес.
    \end{enumerate}
\end{enumerate}

\subsection{Минимальный пример TLB}

Пусть размер страницы $4$ КиБ.

\[
VA = 0x00403010.
\]

Тогда:

\[
VPN = 0x00403,\quad offset = 0x010.
\]

В TLB есть запись:

\[
0x00403 \mapsto 0x001AF.
\]

Следовательно:

\[
PA = 0x001AF000 + 0x010 = 0x001AF010.
\]

\subsection{Page fault}

\begin{definition}
\textbf{Page fault} --- исключение, возникающее при обращении к виртуальной странице, которая отсутствует в оперативной памяти или к которой запрещён требуемый доступ.
\end{definition}

Причины page fault:

\begin{itemize}
    \item страница не загружена в ОЗУ;
    \item обращение к невыделенной памяти;
    \item нарушение прав доступа;
    \item запись в copy-on-write страницу;
    \item попытка исполнения страницы с установленным NX-битом.
\end{itemize}

\subsection{Алгоритм обработки page fault}

\textbf{Вход:} адрес, вызвавший исключение, и тип доступа.

\textbf{Шаги:}

\begin{enumerate}
    \item Процессор сохраняет контекст и передаёт управление обработчику исключения ОС.
    \item ОС определяет виртуальный адрес, вызвавший сбой.
    \item ОС ищет регион памяти, которому принадлежит этот адрес.
    \item Если адрес не принадлежит ни одному допустимому региону:
    \begin{enumerate}
        \item процессу посылается ошибка, например segmentation fault;
        \item обработка завершается.
    \end{enumerate}
    \item Если регион существует, но доступ запрещён:
    \begin{enumerate}
        \item фиксируется нарушение защиты;
        \item процесс обычно завершается.
    \end{enumerate}
    \item Если страница допустима, но отсутствует в ОЗУ:
    \begin{enumerate}
        \item найти свободный кадр;
        \item если свободного кадра нет, выбрать страницу-жертву;
        \item при необходимости записать изменённую страницу на диск;
        \item загрузить требуемую страницу из файла, swap или создать нулевую страницу;
        \item обновить таблицу страниц;
        \item инвалидировать или обновить TLB;
        \item возобновить выполнение инструкции.
    \end{enumerate}
\end{enumerate}

\section{Файлы подкачки и swap}

\subsection{Назначение подкачки}

\begin{definition}
\textbf{Подкачка} --- механизм перемещения страниц между оперативной памятью и внешним хранилищем.
\end{definition}

\begin{definition}
\textbf{Файл подкачки} или \textbf{swap-файл} --- файл на диске, используемый ОС как хранилище выгруженных страниц памяти.
\end{definition}

Также может использоваться отдельный swap-раздел.

Подкачка позволяет:

\begin{itemize}
    \item освободить кадры оперативной памяти;
    \item поддерживать больший объём виртуальной памяти;
    \item временно выгружать редко используемые страницы;
    \item обеспечивать работу demand paging.
\end{itemize}

\subsection{Типы страниц с точки зрения подкачки}

\begin{itemize}
    \item \textbf{Чистые файловые страницы}: содержимое совпадает с файлом на диске; можно просто удалить из ОЗУ.
    \item \textbf{Грязные файловые страницы}: изменены; перед выгрузкой нужно записать обратно в файл.
    \item \textbf{Анонимные чистые страницы}: например нулевые страницы; могут быть восстановлены без чтения с диска.
    \item \textbf{Анонимные грязные страницы}: данные стека или кучи; при вытеснении записываются в swap.
\end{itemize}

\subsection{Алгоритм выгрузки страницы}

\textbf{Вход:} необходимость освободить кадр памяти.

\textbf{Шаги:}

\begin{enumerate}
    \item Выбрать страницу-жертву с помощью алгоритма замещения.
    \item Проверить, закреплена ли страница в памяти.
    \item Если страница закреплена, выбрать другую.
    \item Проверить бит изменения dirty.
    \item Если страница чистая и имеет источник на диске:
    \begin{enumerate}
        \item удалить отображение из таблицы страниц;
        \item пометить кадр свободным.
    \end{enumerate}
    \item Если страница грязная:
    \begin{enumerate}
        \item выделить место в swap или определить файл назначения;
        \item записать страницу на диск;
        \item обновить служебные структуры ОС;
        \item сбросить present bit в PTE;
        \item пометить кадр свободным.
    \end{enumerate}
    \item Инвалидировать соответствующую запись TLB.
\end{enumerate}

\section{Алгоритмы замещения страниц}

\subsection{Постановка задачи}

Если произошёл page fault, а свободных кадров нет, ОС должна выбрать страницу, которую можно выгрузить.

\begin{definition}
\textbf{Алгоритм замещения страниц} --- алгоритм выбора страницы-жертвы для освобождения кадра оперативной памяти.
\end{definition}

Цель алгоритма --- минимизировать число page fault.

\subsection{Оптимальный алгоритм OPT}

\begin{definition}
\textbf{OPT} --- алгоритм, который вытесняет страницу, обращение к которой произойдёт позже всего в будущем.
\end{definition}

\textbf{Шаги алгоритма OPT:}

\begin{enumerate}
    \item При page fault проверить, есть ли свободный кадр.
    \item Если есть, загрузить страницу в свободный кадр.
    \item Если свободного кадра нет:
    \begin{enumerate}
        \item для каждой страницы в памяти определить момент её следующего использования;
        \item выбрать страницу, которая будет использована позже всех;
        \item если страница больше не будет использоваться, выбрать её;
        \item заменить выбранную страницу.
    \end{enumerate}
\end{enumerate}

OPT практически не реализуем в реальной ОС, потому что требует знания будущих обращений, но используется как теоретический эталон.

\subsection{Теорема об оптимальности OPT}

\begin{theorem}
Для заданной последовательности обращений к страницам и фиксированного числа кадров алгоритм OPT порождает минимально возможное число page fault среди всех алгоритмов замещения страниц.
\end{theorem}

\begin{proof}
Докажем обменным аргументом.

Пусть дана последовательность обращений и некоторое число кадров $m$.
Рассмотрим произвольный алгоритм $A$, который отличается от OPT.
Пусть $t$ --- первый момент, когда при необходимости вытеснения алгоритм $A$ выбирает страницу $x$, а OPT выбирает страницу $y$.

По определению OPT страница $y$ используется в будущем не раньше, чем страница $x$.
Возможны два случая.

\textbf{Случай 1.} Страница $y$ больше никогда не используется.
Тогда вытеснение $y$ не может вызвать дополнительных промахов в будущем. Если алгоритм $A$ вместо $x$ вытеснит $y$, он не увеличит число page fault.

\textbf{Случай 2.} Страница $y$ используется позже страницы $x$.
До ближайшего обращения к $y$ страница $x$ понадобится раньше. Поэтому алгоритм, оставивший $x$ в памяти и вытеснивший $y$, не хуже алгоритма, который вытеснил $x$.
В момент, когда потребуется $y$, алгоритм может иметь то же состояние, что и прежний алгоритм, или состояние, отличающееся только перестановкой одной страницы. Число page fault при такой замене не увеличивается.

Следовательно, первый выбор алгоритма $A$, отличающийся от OPT, можно заменить выбором OPT без увеличения числа page fault. Повторяя такое преобразование для всех отличий от OPT, получаем алгоритм, совпадающий с OPT и имеющий не больше page fault, чем $A$.

Так как $A$ был произвольным, OPT минимизирует число page fault.
\end{proof}

\subsection{FIFO}

\begin{definition}
\textbf{FIFO} --- алгоритм, вытесняющий страницу, которая дольше всех находится в памяти.
\end{definition}

\textbf{Шаги FIFO:}

\begin{enumerate}
    \item Поддерживать очередь страниц в памяти.
    \item Новую загруженную страницу помещать в конец очереди.
    \item При необходимости вытеснения выбрать страницу из начала очереди.
    \item Удалить её из памяти.
    \item Загрузить новую страницу и добавить её в конец очереди.
\end{enumerate}

\subsubsection{Пример FIFO}

Пусть есть 3 кадра и последовательность обращений:

\[
1,2,3,4,1,2,5.
\]

\[
\begin{array}{c|c|c}
\text{Обращение} & \text{Содержимое кадров} & \text{Результат} \\
\hline
1 & 1 & fault \\
2 & 1,2 & fault \\
3 & 1,2,3 & fault \\
4 & 2,3,4 & fault,\; вытеснена\;1 \\
1 & 3,4,1 & fault,\; вытеснена\;2 \\
2 & 4,1,2 & fault,\; вытеснена\;3 \\
5 & 1,2,5 & fault,\; вытеснена\;4 \\
\end{array}
\]

FIFO прост, но может работать плохо, потому что не учитывает частоту и давность обращений.

\subsection{Аномалия Белади}

\begin{definition}
\textbf{Аномалия Белади} --- ситуация, когда увеличение числа кадров памяти приводит к увеличению числа page fault для некоторого алгоритма замещения.
\end{definition}

FIFO подвержен аномалии Белади.

Классический пример последовательности:

\[
1,2,3,4,1,2,5,1,2,3,4,5.
\]

Для FIFO с 3 кадрами будет 9 page fault, а с 4 кадрами --- 10 page fault.

\subsection{LRU}

\begin{definition}
\textbf{LRU} --- least recently used, алгоритм, вытесняющий страницу, к которой дольше всего не было обращений.
\end{definition}

\textbf{Шаги LRU:}

\begin{enumerate}
    \item Для каждой страницы хранить время последнего обращения или поддерживать порядок использования.
    \item При каждом обращении к странице обновлять её время.
    \item При page fault и отсутствии свободных кадров выбрать страницу с самым старым временем последнего обращения.
    \item Вытеснить выбранную страницу.
    \item Загрузить новую страницу и отметить её как недавно использованную.
\end{enumerate}

\subsubsection{Пример LRU}

Пусть есть 3 кадра:

\[
1,2,3,4,1,2,5.
\]

\[
\begin{array}{c|c|c}
\text{Обращение} & \text{Кадры} & \text{Результат} \\
\hline
1 & 1 & fault \\
2 & 1,2 & fault \\
3 & 1,2,3 & fault \\
4 & 2,3,4 & fault,\; LRU=1 \\
1 & 3,4,1 & fault,\; LRU=2 \\
2 & 4,1,2 & fault,\; LRU=3 \\
5 & 1,2,5 & fault,\; LRU=4 \\
\end{array}
\]

На этой короткой последовательности LRU совпал с FIFO, но в общем случае LRU обычно лучше учитывает локальность обращений.

\subsection{Теорема о стековом свойстве LRU}

\begin{theorem}
LRU является стековым алгоритмом: множество страниц, находящихся в памяти при $m$ кадрах, всегда является подмножеством множества страниц, находящихся в памяти при $m+1$ кадрах, если обе системы обрабатывают одну и ту же последовательность обращений.
\end{theorem}

\begin{proof}
LRU хранит страницы в порядке убывания давности использования: сверху находится последняя использованная страница, ниже --- менее недавно использованные.

При обращении к странице она переносится наверх стека. Если страница отсутствует, она добавляется наверх. Система с $m$ кадрами хранит первые $m$ элементов этого LRU-стека, а система с $m+1$ кадрами хранит первые $m+1$ элементов.

После каждого обращения порядок LRU-стека определяется только историей обращений, а не числом кадров. Поэтому первые $m$ элементов всегда содержатся среди первых $m+1$ элементов.

Следовательно, множество страниц при $m$ кадрах является подмножеством множества страниц при $m+1$ кадрах.
\end{proof}

\begin{statement}
Из стекового свойства следует, что LRU не подвержен аномалии Белади.
\end{statement}

\begin{proof}
Если при $m$ кадрах требуемая страница находится в памяти, то она принадлежит множеству страниц для $m$ кадров.
По стековому свойству это множество является подмножеством множества страниц для $m+1$ кадров.
Значит, при $m+1$ кадрах эта страница также находится в памяти.
Следовательно, увеличение числа кадров не может превратить попадание в промах.
Общее число page fault не возрастает.
\end{proof}

\subsection{Clock, или алгоритм второго шанса}

\begin{definition}
\textbf{Clock} --- приближённая реализация LRU, использующая бит обращения и циклический список страниц.
\end{definition}

Каждая страница имеет бит обращения $R$.

\textbf{Шаги Clock:}

\begin{enumerate}
    \item Страницы расположены по кругу.
    \item Указатель clock hand показывает на кандидата на вытеснение.
    \item При необходимости вытеснения проверить бит $R$ текущей страницы.
    \item Если $R=0$, выбрать эту страницу для вытеснения.
    \item Если $R=1$:
    \begin{enumerate}
        \item сбросить $R$ в 0;
        \item перейти к следующей странице;
        \item повторить проверку.
    \end{enumerate}
\end{enumerate}

\subsubsection{Пример Clock}

Пусть в памяти страницы:

\[
A,B,C,D.
\]

Биты обращения:

\[
A:1,\quad B:0,\quad C:1,\quad D:0.
\]

Указатель стоит на $A$.

\begin{enumerate}
    \item Проверяем $A$: $R=1$, сбрасываем $R=0$, идём дальше.
    \item Проверяем $B$: $R=0$, выбираем $B$ для вытеснения.
\end{enumerate}

Clock дешевле LRU, потому что не требует точного хранения времени последнего обращения.

\subsection{NRU}

\begin{definition}
\textbf{NRU} --- not recently used, алгоритм, выбирающий страницу из класса, определяемого битами обращения и изменения.
\end{definition}

Используются два бита:

\[
R \text{ --- referenced},\quad M \text{ --- modified}.
\]

Классы:

\[
\begin{array}{c|c|c}
\text{Класс} & R & M \\
\hline
0 & 0 & 0 \\
1 & 0 & 1 \\
2 & 1 & 0 \\
3 & 1 & 1 \\
\end{array}
\]

Алгоритм выбирает страницу из минимального непустого класса.

\textbf{Шаги NRU:}

\begin{enumerate}
    \item Периодически сбрасывать биты $R$.
    \item При необходимости вытеснения разбить страницы на классы по $(R,M)$.
    \item Выбрать случайную или первую страницу из минимального непустого класса.
    \item Вытеснить её.
\end{enumerate}

\subsection{Рабочее множество}

\begin{definition}
\textbf{Рабочее множество процесса} $W(t,\Delta)$ --- множество страниц, к которым процесс обращался в интервале времени от $t-\Delta$ до $t$.
\end{definition}

Формально:

\[
W(t,\Delta)=\{p \mid p \text{ использовалась в интервале } [t-\Delta,t]\}.
\]

Идея: процесс должен иметь в памяти своё рабочее множество. Если суммарные рабочие множества процессов не помещаются в ОЗУ, возникает thrashing.

\begin{definition}
\textbf{Thrashing} --- состояние системы, при котором большая часть времени тратится на подкачку страниц, а не на выполнение полезных вычислений.
\end{definition}

\section{Сегментно-страничная организация}

В некоторых архитектурах возможна комбинация сегментации и страничной организации.

Общая схема:

\[
\text{логический адрес}
=
(\text{segment selector}, \text{offset}).
\]

Сначала:

\[
LA = SegmentBase + Offset.
\]

Затем:

\[
LA \rightarrow PA
\]

через таблицы страниц.

Преимущества комбинации:

\begin{itemize}
    \item сегментация задаёт логические области и права;
    \item paging обеспечивает эффективное размещение в физической памяти;
    \item можно независимо защищать области и страницы.
\end{itemize}

Недостатки:

\begin{itemize}
    \item усложнение аппаратуры;
    \item усложнение ОС;
    \item дополнительные проверки при обращении к памяти.
\end{itemize}

В современных 64-битных системах часто используется почти плоская модель адресации: сегментация минимизирована, а основная защита и виртуализация памяти реализованы через страницы.

\section{Защита памяти}

\subsection{Права доступа}

Память защищается с помощью прав:

\[
R, W, X.
\]

Где:

\begin{itemize}
    \item $R$ --- чтение;
    \item $W$ --- запись;
    \item $X$ --- исполнение.
\end{itemize}

Пример:

\[
\begin{array}{c|c}
\text{Регион} & \text{Права} \\
\hline
Код программы & R,X \\
Глобальные константы & R \\
Данные & R,W \\
Стек & R,W \\
\end{array}
\]

Запрет исполнения данных снижает вероятность эксплуатации переполнений буфера.

\subsection{Copy-on-write}

\begin{definition}
\textbf{Copy-on-write} --- механизм, при котором несколько процессов временно разделяют одну физическую страницу, а копия создаётся только при попытке записи.
\end{definition}

Алгоритм copy-on-write:

\begin{enumerate}
    \item После \texttt{fork} родитель и потомок указывают на одни и те же физические страницы.
    \item Эти страницы помечаются как read-only.
    \item При чтении оба процесса используют общую страницу.
    \item При попытке записи возникает page fault.
    \item ОС проверяет, что это COW-страница.
    \item ОС выделяет новый кадр.
    \item Копирует старую страницу в новый кадр.
    \item Обновляет PTE пишущего процесса.
    \item Разрешает запись и перезапускает инструкцию.
\end{enumerate}

\section{Фрагментация памяти}

\subsection{Внутренняя фрагментация}

\begin{definition}
\textbf{Внутренняя фрагментация} --- неиспользуемая память внутри выделенного блока.
\end{definition}

Например, если размер страницы 4096 байт, а процессу нужна область 4100 байт, потребуется две страницы:

\[
8192 - 4100 = 4092 \text{ байт}
\]

могут остаться неиспользованными.

\subsection{Внешняя фрагментация}

\begin{definition}
\textbf{Внешняя фрагментация} --- наличие свободной памяти, разбитой на мелкие несмежные участки, из-за чего невозможно выделить большой непрерывный блок.
\end{definition}

Сегментация подвержена внешней фрагментации, а чистая страничная организация в основном устраняет её, потому что физические кадры имеют одинаковый размер.

\section{Минимальная модель виртуальной памяти}

Рассмотрим простую систему:

\begin{itemize}
    \item виртуальный адрес: 16 бит;
    \item размер страницы: $256$ байт $=2^8$;
    \item физическая память: 4 кадра;
    \item значит, смещение занимает 8 бит.
\end{itemize}

Виртуальный адрес:

\[
VA = VPN \;|\; offset.
\]

Так как адрес 16-битный:

\[
VPN: 8 \text{ бит},\quad offset: 8 \text{ бит}.
\]

Пусть таблица страниц:

\[
\begin{array}{c|c|c}
VPN & Present & PFN \\
\hline
0x00 & 1 & 0x2 \\
0x01 & 1 & 0x0 \\
0x02 & 0 & - \\
0x03 & 1 & 0x3 \\
\end{array}
\]

Пусть процесс обращается к адресу:

\[
VA = 0x012A.
\]

Тогда:

\[
VPN = 0x01,\quad offset = 0x2A.
\]

По таблице:

\[
VPN=0x01 \mapsto PFN=0x0.
\]

Физический адрес:

\[
PA = 0x00 \cdot 0x100 + 0x2A = 0x002A.
\]

Если бы процесс обратился к адресу:

\[
VA = 0x02F0,
\]

то:

\[
VPN = 0x02,\quad offset = 0xF0.
\]

Но для $VPN=0x02$:

\[
Present = 0.
\]

Следовательно, возник бы page fault.

\section{Краткое сравнение сегментации и paging}

\begin{center}
\begin{tabular}{|p{4cm}|p{5cm}|p{5cm}|}
\hline
\textbf{Критерий} & \textbf{Сегментация} & \textbf{Страничная организация} \\
\hline
Единица разбиения & Логический сегмент переменного размера & Страница фиксированного размера \\
\hline
Отражает структуру программы & Да & Нет, страницы технические \\
\hline
Фрагментация & Внешняя & Внутренняя \\
\hline
Защита & На уровне сегментов & На уровне страниц \\
\hline
Размещение в физической памяти & Обычно требует непрерывной области для сегмента & Страницы могут быть в любых кадрах \\
\hline
Типичное современное применение & Ограниченное, часто плоская модель & Основной механизм виртуальной памяти \\
\hline
\end{tabular}
\end{center}

\section{Итоговая схема преобразования адреса}

В современной системе с виртуальной памятью путь обращения к памяти можно представить так:

\[
\boxed{
\text{Команда процессора}
}
\]

\[
\downarrow
\]

\[
\boxed{
\text{Вычисление эффективного адреса}
}
\]

\[
\downarrow
\]

\[
\boxed{
\text{Логический или виртуальный адрес}
}
\]

\[
\downarrow
\]

\[
\boxed{
\text{Сегментная проверка, если используется}
}
\]

\[
\downarrow
\]

\[
\boxed{
\text{Линейный адрес}
}
\]

\[
\downarrow
\]

\[
\boxed{
\text{TLB lookup}
}
\]

\[
\downarrow
\]

\[
\boxed{
\text{Page table walk при TLB miss}
}
\]

\[
\downarrow
\]

\[
\boxed{
\text{Проверка прав доступа}
}
\]

\[
\downarrow
\]

\[
\boxed{
\text{Физический адрес}
}
\]

\[
\downarrow
\]

\[
\boxed{
\text{Кэш / оперативная память}
}
\]

\section{Основные выводы}

\begin{enumerate}
    \item Оперативная память является центральным уровнем иерархии памяти, но используется совместно с кэшами и внешним хранилищем.
    \item Способы адресации определяют, как процессор вычисляет адрес операнда.
    \item В реальном режиме x86 физический адрес вычисляется как $Segment \cdot 16 + Offset$.
    \item В защищённом режиме используются дескрипторы сегментов, права доступа, уровни привилегий и paging.
    \item Виртуальная память предоставляет каждому процессу собственное адресное пространство.
    \item Страничная организация отображает виртуальные страницы в физические кадры.
    \item MMU и TLB аппаратно ускоряют трансляцию адресов.
    \item Page fault не всегда является ошибкой: часто это штатный механизм demand paging.
    \item Файл подкачки хранит страницы, временно выгруженные из оперативной памяти.
    \item Алгоритмы замещения страниц выбирают страницу-жертву при нехватке кадров.
    \item OPT теоретически минимален, LRU использует локальность, Clock является практичным приближением LRU.
    \item Сегменты, страницы и регионы описывают память на разных уровнях: логическом, аппаратном и уровне ОС.
\end{enumerate}

\end{document}


\end{document}
